\subsection{A polynomial causal inverse in the gain norm}
\label{sec:short-polynomial-causal}

Retain the fixed spaces, actual train, and gain from
Proposition~\ref{prop:short-full-clock}. Put
\[
 \mathcal H=L^2(s^{-3/2}ds;L^2_{\rm div}(dq/c\,d\eta)),\quad
 \mathcal T=L^2(s^{-1/2}ds;L^2(dq/c\,d\eta)).
\]
Let $K$ be a fixed compact subset of $\{0<\Re\mu<1/4\}$ and
$k_*:=\inf_K(\Re\mu+1/4)>0$. Constants below are uniform on $K$.

\begin{proposition}\label{prop:short-polynomial-causal}
For every planar symmetric tensor $F:C_c\times(0,\infty)\to
\operatorname{Sym}_2$ with $g^{-1}F\in\mathcal T$, there is
a unique causal zero-entrance energy solution of
\begin{equation}\label{eq:short-causal-equation}
 b_s+\frac\mu s b-\nu\Delta b+\mathbb P_0\operatorname{div}
       (v\otimes b+b\otimes v)=\mathbb P_0\operatorname{div}F.
\end{equation}
Here zero entrance means $b\in\mathcal H$ near zero; no homogeneous
$s^{-\mu}$ entrance is adjoined. The energy class also requires
$\nabla b\in L^2_{\rm loc}((0,\infty);L^2)$. There is a finite
$M$ such that
\begin{equation}\label{eq:short-causal-bound}
 \|g^{-1}b\|_{\mathcal H}
 +\sqrt\nu\|\nabla(g^{-1}b)\|_{\mathcal T}
 \le C\nu^{-M}\|g^{-1}F\|_{\mathcal T}.
\end{equation}
The response is locally continuous in $\mu$ with the geometry fixed.
This is a bound after conjugation by $g$; its ordinary norm can contain
$e^{\gamma/\nu}$.
\end{proposition}

\begin{proof}
Use Proposition~\ref{prop:short-full-clock} with its fixed $b$,
$H=\nu^{-1}((\nu b)^{-1}-1)$, and weak-source degree $M_s$.

\emph{The zero-entrance construction and all ages.}
Put $T=g^{-1}F$, $u=g^{-1}b$ and
$h=g'/g=\gamma/(1+\nu s)^2\ge0$. Solve the free equation for $u_0$
with scalar term $\mu/s+h$ and tensor $T$ by starting from zero at
$s=\delta$ and taking $\delta\downarrow0$. The weighted energy is
coercive and gives
\[
 \|u_0\|_{\mathcal H}\le(2\sqrt{k_*\nu})^{-1}\|T\|_{\mathcal T},
 \qquad \|\nabla u_0\|_{\mathcal T}\le\nu^{-1}\|T\|_{\mathcal T}.
\]
Spatial coefficient estimates may first be taken with unnormalized
cell measure. Multiplying velocity, tensor, $X_m$, and weighted-dual
source norms by the same $c^{-1/2}$ gives the normalized convention
used in $\mathcal H$ and $\mathcal T$; all inequalities below are
unchanged. Equivalently define the source duality using normalized
cell pairing and normalized test norm.
The localized vorticity source
$f=-\operatorname{curl}(v\cdot\nabla u_0+u_0\cdot\nabla v)$
has zero mass and satisfies
\[
 \|f(s)\|_{Z^{-1}_{m,c}}
 \le A_m(\|\nabla u_0(s)\|_2+\|u_0(s)\|_2),\qquad
 A_m\le C\nu^{-p_m}\quad(0<s<H),
\]
for one fixed finite $p_m$. Indeed $A_m$ is bounded by a constant
times $\sup_{s\le H}\|\langle x\rangle^m(v,\nabla v)(s)\|_\infty$.
The actual train tail \eqref{sct:transverse}, with $S_*=H$, bounds
this by a fixed polynomial in $c,\nu^{-1},H$: the nonzero modes
have exponential transverse kernels, and the zero mode has decaying
primitives at both ends. Here $c=\nu^{-2}$ and $H=O(\nu^{-2})$.
Thus this is a polynomial coefficient bound, with no exponential
weighted-energy loss. Testing the outer curl against $H^1_{-m}$
then gives the displayed negative-order estimate. Hence
$\|f\|_{L^2(s^{-1/2};Z^{-1}_m)}
\le C\nu^{-p_m-3/2}\|T\|_{\mathcal T}$.

For $\sigma=\Re\mu$, the substitution $t=su$ and Minkowski give
the exact Hardy inequality
\begin{equation}\label{eq:short-causal-hardy}
 \left\|\int_0^s(t/s)^\sigma(s-t)^{-1/2}f(t)dt\right\|_{L^2(s^{-3/2})}
 \le B(\sigma+3/4,1/2)\|f\|_{L^2(s^{-1/2})}.
\end{equation}
The kernel $1$ on $s<H$ costs an additional $\sqrt H\le C_b\nu^{-1}$.
Define using only forward propagation
\[
 \omega(s)=\int_0^s(t/s)^\mu\frac{g(t)}{g(s)}U(s,t)f(t)dt,
 \qquad u=u_0+K_c\omega.
\]
Equations \eqref{eq:short-causal-weak-forward}--\eqref{eq:short-causal-hardy}
give $\|\omega\|_{L^2(s^{-3/2};X_m)}
\le C\nu^{-M_s-p_m-5/2}\|T\|_{\mathcal T}$.
Zero-circulation Biot--Savart reconstruction gives
$\|K_c\omega\|_{H^1}\le C_m(1+c)\|\omega\|_{X_m}$:
nonzero horizontal modes cost $c$, while the mean is the
zero-endpoint primitive controlled by one-dimensional Hardy:
$\|\int_\eta^\infty\bar\omega(\zeta)d\zeta\|_2
\le2\|\eta\bar\omega(\eta)\|_2$. This follows by integrating
$(\eta|\int_\eta^\infty\bar\omega|^2)'$; zero circulation
removes both endpoint terms.
Since $\sqrt H=O(\nu^{-1})$, we obtain on the core
\[
 \|u\|_{\mathcal H(0,H)}
 +\sqrt\nu\|\nabla u\|_{\mathcal T(0,H)}
 \le C\nu^{-M_s-p_m-5}\|T\|_{\mathcal T(0,H)}.
\]
The velocity equation follows from its curl and divergence equations:
after testing in age the discrepancy is spatial $H^{-1}$, curl-free
and divergence-free; its Fourier transform is supported only at zero
and must vanish. This retains the horizontal mean. Local Hilbert
triple regularity gives a strong $L^2$ trace at every positive age.

For $E=s^{-1/2}\|u\|_2^2$, $L=\|\operatorname{sym}\nabla v\|_\infty$
and $A(s)=\exp(-2\int_{t_*}^sL)$, the exact energy inequality is
\begin{equation}\label{eq:short-causal-tail-energy}
 (AE)'+2k_*AE/s+2hAE+\nu A s^{-1/2}\|\nabla u\|_2^2
 \le\nu^{-1}A s^{-1/2}\|T\|_2^2.
\end{equation}
Selecting a time in $[H-1,H-1/2]$ and applying this inequality to $H$
gives
$H^{-1/4}\|u(H)\|_2
\le C_b[\sqrt H\|u\|_{\mathcal H(H-1,H)}
+\nu^{-1/2}\|T\|_{\mathcal T(H-1,H)}]$.
With $a_b=1+\nu H=(\nu b)^{-1}$, the actual train estimate
\eqref{short-bg:train-late} yields
\[
 I_H:=\int_H^\infty L(s)ds
 \le C/(\nu a_b)+Ce_c=Cb+C\nu^3.
\]
Thus continuation from the exact strong trace, with the entire
prescribed tensor after $H$, costs only the fixed factor $e^{Cb+C\nu^3}$.
It crosses the overlap regime and reaches infinite age without a
vorticity moment assumption on the late response. Combining the
core and tail gives \eqref{eq:short-causal-bound}, for example with
any $M\ge M_s+p_m+6$.

For uniqueness, a homogeneous difference in the energy class has
a sequence $\delta_j\downarrow0$ with
$\delta_j^{-1/2}\|u(\delta_j)\|_2^2\to0$, because its integral
against $ds/s$ is finite. Apply \eqref{eq:short-causal-tail-energy}
on $[\delta_j,B]$ and let $j\to\infty$. This also proves causality.
Continuity in $\mu$ on the core follows from the same Hardy estimate
with one factor $|\log(t/s)|$; its beta integral is finite uniformly
on $K$. In the free and tail equations the difference contributes
$-(\mu-\mu')u_{\mu'}/s$, whose weighted pairing is bounded by
$|\mu-\mu'|\|u_{\mu'}\|_{\mathcal H}\|u_\mu-u_{\mu'}\|_{\mathcal H}$
and is absorbed by the $k_*$ coercivity. The local energy estimate
also gives trace continuity. All geometric parameters stay fixed.
\end{proof}

\begin{corollary}\label{cor:short-polynomial-stress}
Let $S_{0,\mu}$ and $S_{\rm lin,\mu}$ denote the free and actual
causal tensor responses, acting on the sum of the two components of
the split stress space $\mathcal Z=\mathcal Z_1\oplus
\rho_b^{-1/2}\mathcal T$, with its specified norm and $0<\rho_b\le1$.
For a three-component cylindrical tensor these flat maps act on its
$(r,z)$ block; its hoop component is retained in the exact curved
comparison.
The first-component weight is at least one, so
$\|F_1+F_2\|_{\mathcal T}\le\|F_1\|_{\mathcal Z_1}
+\sqrt{\rho_b}\|F_2\|_{\rho_b^{-1/2}\mathcal T}$. Let
$M_0u=(v\otimes u+u\otimes v,0)$ and suppose its actual train
multiplier bound is $\|M_0u\|_{\mathcal Z_1}
\le C\nu^{-p_V}\|u\|_{\mathcal H}$. Then
\[
 D_0^{-1}=(I+M_0S_{0,\mu})^{-1}=I-M_0S_{\rm lin,\mu},\qquad
 \|g^{-1}D_0^{-1}g\|_{\mathcal Z\to\mathcal Z}\le C\nu^{-M},
 \quad M=\lceil p_m+p_V+M_s+10\rceil.
\]
Indeed $b=S_{\rm lin,\mu}F$ satisfies
$b=S_{0,\mu}(F-M_0b)$ by free uniqueness, giving both inverse
identities. The second split input component is retained. The degree
has slack for the fixed norm equivalences and trace estimate, and
depends only on the fixed weights and coefficient budgets.
\end{corollary}
